\section{The joint coefficient algebra and observable conditioning}
\label{sec:joint}
The affine residues have a polynomial counterpart, but residues at
individual roots are not the appropriate first decomposition. Distinct
component polynomials may share a root. We separate the polynomials by
their complete coefficient tuples before estimating their scalar zeros.
This distinction is what preserves the additive normalization loss.

Let $\mathcal R_d$ be the model image on which $\tau(P)>0$ and
$\rank P(T_1)=k$. Recover $q$ and $K$ by the normalization equation and
interpolation. The leading coefficient is
$A=K_d=U\diag(\alpha)V^{\mathsf T}$. Choose orthonormal frames
$Q_1,Q_2$ for its column and row spaces, and put
\[
 M(z)=Q_1^{\mathsf T}K(z)Q_2=\sum_{j=0}^dM_jz^j,
 \qquad A_0=M_d,\qquad C_j=M_jA_0^{-1}\quad(j<d).
\]
These matrices are determined up to orthogonal changes of frames. If
$S=Q_1^{\mathsf T}U$ and $T=Q_2^{\mathsf T}V$, then
\begin{equation}\label{eq:joint-algebra}
 C_j=S\diag(f_{bj})S^{-1},\qquad
 M(z)A_0^{-1}=S\diag(f_b(z))S^{-1}.
\end{equation}
The $C_j$ therefore commute and generate a semisimple real algebra.
For each distinct monic polynomial $f$ occurring as a joint eigenvalue
tuple, let $E_f$ denote its joint spectral projector. Concretely, choose
real $c_j$ for which the numbers $\sum_{j<d}c_jf_j$ are distinct for
these finitely many tuples, and take the spectral projectors of
$\sum_{j<d}c_jC_j$. The forbidden choices form finitely many proper
hyperplanes. The resulting projectors are independent of this choice.
Set
\begin{equation}\label{eq:joint-condition}
 W_f=A_0^{-1}E_f,\qquad
 B_d(P)=\sum_{f\ {\rm distinct}}\|W_f\|_\op,
 \qquad \eta_d(P)=(B_d(P)+\tau(P)^{-1})^{-1}.
\end{equation}
Set $\eta_d=0$ off $\mathcal R_d$. This is a finite construction from
the data. No optimization over latent factorizations enters it. The
rank of $E_f$ records the number of components with polynomial $f$.

\begin{theorem}[Observable polynomial inverse]\label{thm:joint}
The quantities in \eqref{eq:joint-condition} are independent of the
orthonormal frames and of the latent representation. On $\mathcal R_d$,
\begin{equation}\label{eq:joint-comparison}
 \|A_0^{-1}\|\le B_d(P)\le \frac{k}{\alpha_*\kappa},\qquad
 \eta_d(P)\ge c(\kappa^{-1}+\tau(P)^{-1})^{-1}.
\end{equation}
Every competitor in the entire closed model $\cK_d$ satisfies
\begin{equation}\label{eq:joint-inverse}
 d_\infty(\Lambda,\Lambda')
 \le C\min\{1,[\delta/\eta_d(P)]^{1/d}\}.
\end{equation}
If the distinct roots within each true component polynomial have
separation at least $\nu>0$ and multiplicity at most $m$, the exponent
is $1/m$, with an additional dependence on $\nu$. In particular,
within-component simple roots give a linear estimate through arbitrary
cross-component collisions. The same estimates hold for the direct
estimator of Section~\ref{sec:algorithm}, with its error $e$ in place
of $\delta$ and no conditioning floor supplied to the estimator.

On $\mathcal R_d$, $B_d$ is lower semicontinuous and is continuous on
strata with a fixed partition into equal component polynomials.
The zero extension of $\eta_d$ is upper semicontinuous on the whole
model image, and
\[
 \eta_d(P)>0\quad\Longleftrightarrow\quad
 \tau(P)>0\ \hbox{and}\ \rank P(T_1)=k.
\]
For $d=1$, $B_1=\beta$ and $\eta_1\asymp\eta$, with the notation of
Section~\ref{sec:residue}.
\end{theorem}

\begin{proof}
Formula \eqref{eq:joint-algebra} gives
\begin{equation}\label{eq:resolved-inverse}
 E_f=S\diag(\mathbf1_{f_b=f})S^{-1},\qquad
 W_f=T^{-\mathsf T}\diag(\alpha_b^{-1}\mathbf1_{f_b=f})S^{-1},
 \qquad M(z)^{-1}=\sum_f\frac{W_f}{f(z)}.
\end{equation}
The algebra, its joint eigenvalue tuples and its joint projectors are
observable. An orthogonal change $Q_1\mapsto Q_1O_1$,
$Q_2\mapsto Q_2O_2$ sends $W_f$ to $O_2^{\mathsf T}W_fO_1$.
The norm is unchanged. Since $\sum_fW_f=A_0^{-1}$, the triangle
inequality gives the lower bound in \eqref{eq:joint-comparison}.
Each summand is bounded by $(\alpha_*\kappa)^{-1}$; there are at most
$k$ of them. This proves both comparisons.

Write $r=q'-q$. The normalization identity
\eqref{eq:normalization-difference} gives
$\|r\|_{\rm coef}\le C\delta/\tau$. Here and below the scalar
coefficient norm is Euclidean or the sum norm; their fixed dimensional
comparison is included in $C$. Let $I_d$ be interpolation at the first
$d+1$ clocks and define the scalar polynomial
\[
 h_f[r](z)=I_d\big(r(T_j)f(T_j)/q(T_j)\big)(z).
\]
It is linear in $r$, and $\|h_f[r]\|_{\rm coef}\le C\|r\|_{\rm coef}$
uniformly over the model. The bounds use only clock geometry and the
positive bounds for $f(T_j)/q(T_j)$, not channel inverses. Form
\begin{equation}\label{eq:polynomial-gauge}
 \overline M_r(z)=S\diag\big(\alpha_b(f_b(z)+h_{f_b}[r](z))\big)T^{\mathsf T}.
\end{equation}
This is the reduction of $K+I_d(rP)$. It is used algebraically; its
scalar roots need not remain real. It satisfies the exact identity
\begin{equation}\label{eq:gauge-resolved}
 \overline M_r(z)^{-1}=\sum_f\frac{W_f}{f(z)+h_f[r](z)}.
\end{equation}
In particular equal complete polynomials undergo the same gauge. There
is no need to split a residue at a root shared by different polynomials.

The reduction of the competitor has the form
$Q_1^{\mathsf T}K'(z)Q_2=\overline M_r(z)+H(z)$, where $H$ interpolates
$q'(T_j)Q_1^{\mathsf T}(P'_j-P_j)Q_2$. Hence
$\|H\|_{\rm coef}\le C\delta$. On a fixed disk containing all roots,
put $\rho=\dist(z,\Lambda)$. Since $|f(z)|\ge\rho^d$, whenever
$\rho^d\ge C_1\|r\|_{\rm coef}$ the scalar denominators in
\eqref{eq:gauge-resolved} have modulus at least $\rho^d/2$. Thus
\[
 \|\overline M_r(z)^{-1}H(z)\|
 \le C B_d\delta\rho^{-d}.
\]
Choose $a^d=C_2(\|r\|_{\rm coef}+B_d\delta)$, with $C_2$ large.
Every determinant on the homotopy $\overline M_r+sH$ is nonzero outside
the $a$-disks about the true roots. The same is true on the preliminary
homotopy $M+u(\overline M_r-M)$, since its scalar denominators obey the
same estimate. Leading coefficients stay invertible: the scalar leading
coefficients in \eqref{eq:polynomial-gauge} are close to one, and their
reduced inverse has norm at most $2B_d$; then $H_d$ is a small relative
perturbation. For large $|z|$, $|f(z)|\ge c|z|^d$ and both errors are
bounded by their coefficient norms times $C|z|^d$. The same Neumann
argument excludes roots there. There are therefore exactly $kd$ roots
through both homotopies. The contour and fixed-permutation limiting
argument of Lemma~\ref{lem:roots} gives the matching in
\eqref{eq:joint-inverse}. An invertible leading coefficient of the
reduced competitor forces its two reduced channel factors to be
invertible, so these are exactly its $kd$ component roots. Larger
errors are covered by the bounded target interval.

For the refined exponent, near a root cluster of a given $f$, all
factors outside that cluster are bounded below in terms of $\nu$.
There are at most $m$ factors in the cluster. On a fixed disk this gives
$|f(z)|\ge c_\nu\min\{1,\rho^m\}$. Repeat the two homotopies with
$a^m=C_\nu(\|r\|_{\rm coef}+B_d\delta)$. This argument is also
uniform on a fixed system of separated clusters containing at most
$m$ roots of each component: gaps inside the clusters may vanish.

For the estimator, $\|\widehat q-q\|_{\rm coef}\le Ce/\tau$ as in
\eqref{eq:algorithm-q}. Formula \eqref{eq:resolved-inverse} at $T_1$
gives $\sigma_k(P_1)^{-1}\le CB_d$. The estimated row and column
frames therefore have invertible overlaps with the true frames,
with inverse norms at most two, when $eB_d$ is small. In those frames
each $W_f$ is multiplied on the left and right by these inverse
overlaps, and their norm sum is at most $4B_d$. The preceding proof
applies verbatim with $r=\widehat q-q$ and interpolation error $Ce$.
All fallback conditions are then excluded. Projecting complex roots
to real parts and to $[L,D]$ cannot increase their matched distances.

For semicontinuity, take a sequence of data converging inside
$\mathcal R_d$ and choose continuous local orthonormal frames. The
normalizations, coefficients and the matrices $C_j$ converge. Choose
one linear combination separating the distinct joint tuples of the
limit. Its contour projectors give, for each limiting polynomial $f$,
\[
 \sum_{g\ \mathrm{in\ the\ cluster\ of}\ f} W_g^{(n)}\longrightarrow W_f.
\]
To justify clustering of the full tuples, one may also take a compact
subsequence of the original model parameters. The rank-$k$ limiting
observation forces both limiting channels to have full rank; their
component coefficients consequently give exactly those joint tuples.
The triangle inequality proves lower semicontinuity. If groups neither
split nor merge, every summand converges, giving continuity. Since
$\tau$ is continuous and positive here, the assertions for $\eta_d$
follow locally. At a limit with $\tau=0$, use $\eta_d\le\tau$.
At a limit with $\tau>0$ and rank loss, the normalizations converge and
\eqref{eq:resolved-inverse} gives
$B_d\ge c\sigma_k(P_1)^{-1}\to\infty$. This proves the zero-extension
statement. Finally, in degree one the distinct polynomials are precisely
$z-x$, so $W_{z-x}$ is the residue $R_x$ and $\tau\asymp\gamma$.
\end{proof}

\subsection{Explicit dependence on degree and clock geometry}
The following estimate makes the constants in the uniform polynomial
bound quantitative. It is not a claim that the monomial interpolation
problem stays well conditioned as degree grows.
Let $\ell_j$ be the Lagrange polynomials at $T_1,\ldots,T_{d+1}$ and put
\begin{align*}
 R_0&=\max\{1,|L|,|D|\},& R&=2R_0+2,&
 \mathcal L&=\sum_{j=1}^{d+1}\|\ell_j\|_{\rm coef},\\
 Q&=\max_{j\le\ell}(T_j-L)^d,&
 V_d&=\max_{j\le\ell}\Big(\sum_{i=0}^{d-1}|T_j|^{2i}\Big)^{1/2},&
 \chi&=\max_{j\le\ell}\left(\frac{T_j-L}{T_j-D}\right)^d,\\
 \mathcal A(P)&=\mathcal L Q\left(B_d+\chi V_d\sqrt\ell/\tau\right).
\end{align*}
Here the scalar coefficient norm is the sum of absolute coefficients.
\begin{proposition}[A quantitative bound]\label{prop:explicit}
If $\delta\mathcal A(P)\le2^{-d-4}$, then
\begin{equation}\label{eq:explicit}
 d_\infty(\Lambda,\Lambda')
 \le (2kd-1)\big(4R^d\delta\mathcal A(P)\big)^{1/d}.
\end{equation}
There is no hidden dependence on $m_1,m_2$ in this deterministic formula:
the data norm is Frobenius, and $B_d,\tau$ already use that observation.
\end{proposition}
\begin{proof}
The Euclidean coefficient norm of $r$ is at most
$\sqrt\ell Q\delta/\tau$. Hence
$\|h_f[r]\|_{\rm coef}\le\mathcal L\chi V_d\sqrt\ell Q\delta/\tau$,
whereas $\|H\|_{\rm coef}\le\mathcal L Q\delta$.
Their sum after weighting the latter by $B_d$ is at most
$\delta\mathcal A$. On $|z|\le R$, evaluation costs at most $R^d$.
With $a^d=4R^d\delta\mathcal A$, both homotopy exclusions above are
strict outside the $a$-disks. On $|z|\ge R$ one has
$|f(z)|\ge2^{-d}|z|^d$; the imposed smallness makes the scalar and
matrix relative errors less than $1/4$. It also ensures invertibility
of every leading coefficient. A disk component with $n$ centres,
counted with multiplicity, admits matching distance at most $(2n-1)a$.
Use $n\le kd$ and pass to nonexceptional radii as in Lemma~\ref{lem:roots}.
\end{proof}
